\chapter{Considerações Finais}
\label{cap:5}

Este capítulo sintetiza os resultados do desenvolvimento e as condições de interpretação das evidências.

\section{Resultados em Relação aos Objetivos}
\label{sec:5.1}

O Quadro\nobreakspace{}\ref{qua:15} relaciona as evidências disponíveis aos cinco objetivos da seção\nobreakspace{}\ref{sec:1.4}.

\begin{TCTabela}{quadro}{qua:15}{Situação dos objetivos específicos}
  {\dimexpr\linewidth-6\tabcolsep-4\arrayrulewidth\relax}{3}{|>{\raggedright\arraybackslash}p{0.35\TTCellWidth}|>{\raggedright\arraybackslash}p{0.15\TTCellWidth}|>{\raggedright\arraybackslash}p{0.5\TTCellWidth}|}
  {\textbf{Objetivo} & \textbf{Etapa} & \textbf{Situação e evidência}}
  {elaboração dos autores.}
1 --- Delimitar e projetar & TCC 1 & Atendido. Lean Inception, jornadas e piloto formativo com oito sessões orientaram a concepção e identificaram questões de linguagem e navegação (seções\nobreakspace{}\ref{sec:3.3} e \ref{sec:4.14}). \\ \hline
2 --- Consolidar a implementação A/B & TCC 2 & Implementado na versão arquivada. Duração ativa, oferta, exposição e resposta alimentam o painel por fluxo (seções\nobreakspace{}\ref{sec:4.11} e \ref{sec:4.12}). \\ \hline
3 --- Registrar e apresentar indicadores & TCC 1 & Atendido no escopo do A/B. O formulário e o painel documentam satisfação, atributos e facilidade com número de respostas válidas. O denominador por fluxo faz parte do objetivo 2 (seção\nobreakspace{}\ref{sec:4.11}). \\ \hline
4 --- Verificar e demonstrar a integração & TCC 2 & Com evidência inicial. As duas jornadas foram demonstradas até o painel. A suíte Python teve cobertura medida e o relógio passou em teste JavaScript (seções\nobreakspace{}\ref{sec:3.6.3} e \ref{sec:4.19}). \\ \hline
5 --- Revisar e verificar as regras & TCC 2 & Parcial. Simulação reproduzida e sensibilidade calculada por perfil. Testes de contratos de acesso e entrada executados. Carga e acessibilidade tiveram medições iniciais. A revisão e a comparação em cenários reservados permanecem no plano (seções\nobreakspace{}\ref{sec:3.6.3}, \ref{sec:4.9} e \ref{sec:4.19}). \\* \hline
\end{TCTabela}

\section{Limitações}
\label{sec:5.2}

O piloto formativo teve quatro participantes por fluxo, recrutamento voluntário, alocação descrita como sorteio sem registro auditável e popup de redirecionamento em ambos os fluxos. Familiaridade, intervenções de moderação e diferenças de exposição restringem interpretações comparativas. O fluxo A carece de registro individual de exposição ao formulário final. Os comentários preservados vieram do questionário pós-sessão, e o piloto carece de um corpus de comentários do formulário para avaliar especificidade.

O piloto adotou TCLE, pseudônimos e ausência de áudio ou câmera. A ausência de apreciação por Comitê de Ética em Pesquisa motivou o encerramento da coleta no piloto. O item de facilidade com a redação da seção\nobreakspace{}\ref{sec:2.7} ainda não foi testado com pessoas. Os protocolos com cidadãos, gestores e especialistas são prospectivos.

A simulação descreve como as regras operam sob suposições definidas pelos autores e não demonstra detecção de dificuldade real. Os testes atuais dizem respeito a contratos selecionados da versão local. A adequação das regras em serviços reais, a utilidade administrativa e a conformidade integral de segurança e acessibilidade exigem investigações próprias.

\section{Trabalhos Futuros}
\label{sec:5.3}

Trabalhos futuros poderão executar os protocolos do Apêndice\nobreakspace{}\ref{ap:K} após a apreciação ética aplicável. H1 propõe que o convite contextual produza maior taxa de resposta que o formulário exibido ao final. H2 propõe que os comentários obtidos por convite contextual sejam mais específicos sobre a etapa vivenciada. O protocolo deverá fixar denominadores, critérios de exposição, rubrica e tratamento da não resposta antes da coleta. A avaliação poderá incluir a interpretação do painel por gestores e inspeção independente por especialistas.

\section{Conclusão}
\label{sec:5.4}

O protótipo A/B associa eventos de navegação a feedback sobre a etapa do serviço. O piloto formativo no Figma orientou a concepção da jornada. Em ambiente local, a verificação técnica percorreu os dois fluxos, confirmou a persistência de exposições e respostas e conciliou os denominadores com o painel.

A regra inicial apresentou discriminação próxima do acaso nas condições sintéticas. A análise por gatilho identificou omissões recuperadas por retornos e acionamentos do leitor lento associados ao tempo ativo. A entrega central do TCC 2 é fundamentar perfis com sinais sobrepostos, revisar a regra e medir os resultados em cenários reservados.
